\section{量子電磁気学}

\noindent\textbf{この章の中心命題}
\[
\boxed{
\mathcal{L}_{\mathrm{QED}}
=
\bar{\psi}(i\gamma^\mu D_\mu-m)\psi
-\frac{1}{4}F_{\mu\nu}F^{\mu\nu},
\qquad
D_\mu=\partial_\mu+ieA_\mu
}
\]
局所 $U(1)$ ゲージ対称性を要請すると、電子場の微分は共変微分に置き換わり、相互作用項 $-e\bar{\psi}\gamma^\mu A_\mu\psi$ が定まる。

\subsection{自由な電子場と電磁場}
量子電磁気学とは、電子・陽電子を表すディラック場 $\psi(x)$ と、電磁相互作用を表す電磁場 $A_\mu(x)$ を同時に扱う理論である。ここで添字 $\mu=0,1,2,3$ は、時間成分一つと空間成分三つをまとめて表している。論理の出発点は、まず自由理論を書き、その後で相互作用を許す条件を課すことである。

量子化後の状態空間は、電子・陽電子のフェルミフォック空間と光子のボースフォック空間のテンソル積
\[
\mathcal{H}_{\mathrm{QED}}
=
\mathcal{F}_-(\mathfrak{h}_e\oplus \mathfrak{h}_{\bar e})
\otimes
\mathcal{F}_+(\mathfrak{h}_\gamma)
\]
である。本ノートでは光子の一粒子空間を
\[
\mathfrak{h}_\gamma
:=
L^2\!\left(\mathbb{R}^3,\frac{d^3p}{(2\pi)^3};\mathbb{C}^2\right)
\]
と書く。ここで $\mathbb{C}^2$ は、光子がもつ二つの独立な物理的内部自由度を表している。したがって QED では、電子・陽電子・光子の自由度をすべてまとめた大きなヒルベルト空間の上で理論を立てることになる。

まず、相互作用を入れる前の自由な場を考える。ディラック場のラグランジアン密度は
\[
\mathcal{L}_{\mathrm{Dirac}}
=
\bar{\psi}(i\gamma^\mu \partial_\mu - m)\psi
\]
である。ここで
\[
\bar{\psi}:=\psi^\dagger \gamma^0
\]
はディラック共役である。一方、電磁場は4成分のポテンシャル $A_\mu(x)$ で表し、その導関数から
\[
F_{\mu\nu}:=\partial_\mu A_\nu-\partial_\nu A_\mu
\]
を作る。$F_{\mu\nu}$ は電場と磁場をひとまとめにした量であり、これを用いると自由な電磁場のラグランジアン密度は
\[
\mathcal{L}_{\mathrm{EM}}
=
-\frac{1}{4}F_{\mu\nu}F^{\mu\nu},
\]
と書ける。

したがって、相互作用を入れる前の全ラグランジアンは
\[
\mathcal{L}_0
=
\bar{\psi}(i\gamma^\mu \partial_\mu - m)\psi
- \frac{1}{4}F_{\mu\nu}F^{\mu\nu}
\]
となる。この段階では、電子場と電磁場はまだ独立に存在している。

\subsection{一様な位相変換と電荷保存}
複素場 $\psi(x)$ 全体に、絶対値1の複素数 $e^{-i\alpha}$ を掛ける変換
\[
\psi(x)\mapsto e^{-i\alpha}\psi(x),
\qquad
\bar{\psi}(x)\mapsto \bar{\psi}(x)e^{i\alpha}
\]
を考える。ここで $\alpha$ は定数であり、全ての時空点で同じ値をとる。このような変換を一様な位相変換という。絶対値1の複素数全体はしばしば $U(1)$ と書かれるので、この対称性を $U(1)$ 対称性と呼ぶ。

ディラック場の自由ラグランジアンはこの変換の下で不変である。したがってネーターの定理により、対応する保存電流
\[
j^\mu = \bar{\psi}\gamma^\mu \psi,
\qquad
\partial_\mu j^\mu = 0
\]
が得られる。時間成分 $j^0$ は電荷密度、空間成分 $\mathbf{j}$ は電流密度であり、この保存則は電荷保存を表している。

\subsection{時空点ごとに位相を変えると何が起こるか}
次に、この位相変換を各時空点ごとに別々に行いたいとする。すなわち、$\alpha$ を定数ではなく関数 $\alpha(x)$ に拡張し、
\[
\psi(x)\mapsto e^{-i\alpha(x)}\psi(x)
\]
を考える。全ての点で同じ位相を掛けるのではなく、点ごとに異なる位相を掛けるのである。このような変換を局所的な位相変換という。

ところが、通常の微分 $\partial_\mu \psi$ は
\[
\partial_\mu \psi \mapsto e^{-i\alpha(x)}
\left(\partial_\mu \psi - i(\partial_\mu \alpha)\psi\right)
\]
と変換し、右辺に $-(\partial_\mu\alpha)\psi$ という余分な項が現れる。したがって、自由ラグランジアンのままでは局所的な位相変換に対して不変でなくなる。

\subsection{電磁場の導入と共変微分}
この余分な項を打ち消すために、新しい場 $A_\mu(x)$ を導入し、
\[
D_\mu := \partial_\mu + ieA_\mu
\]
で定まる微分 $D_\mu$ を考える。これは、局所位相変換で生じる余分な項を $A_\mu$ の変換で打ち消すように定義した微分であり、共変微分と呼ばれる。ここで $e$ は電子の電荷の大きさを表す定数である。

同時に、$A_\mu$ 自身も
\[
A_\mu \mapsto A_\mu - \frac{1}{e}\partial_\mu \alpha
\]
と変換すると、
\[
D_\mu \psi \mapsto e^{-i\alpha(x)}D_\mu \psi
\]
が成り立つ。つまり、$D_\mu\psi$ は $\psi$ と同じ仕方で変換する。これにより、局所的な位相変換の下でも理論を不変にできる。

このように、時空点ごとの位相の取り方を変えても理論の形が変わらないという性質を、局所 $U(1)$ ゲージ対称性と呼ぶ。ここで「ゲージ」とは、記述の仕方の違いにすぎず、物理的内容を変えない自由度のことである。

こうして得られる量子電磁気学のラグランジアンは
\[
\mathcal{L}_{\mathrm{QED}}
=
\bar{\psi}(i\gamma^\mu D_\mu - m)\psi
- \frac{1}{4}F_{\mu\nu}F^{\mu\nu}
\]
である。これを展開すると
\[
\mathcal{L}_{\mathrm{QED}}
=
\bar{\psi}(i\gamma^\mu \partial_\mu - m)\psi
- \frac{1}{4}F_{\mu\nu}F^{\mu\nu}
- e\bar{\psi}\gamma^\mu A_\mu \psi
\]
となり、相互作用項
\[
\mathcal{L}_{\mathrm{int}} = - e\bar{\psi}\gamma^\mu A_\mu \psi = -e j^\mu A_\mu
\]
が現れる。すなわち、電子の電流 $j^\mu$ が電磁場 $A_\mu$ と結合することで相互作用が生じる。重要なのは、この形が思いつきで置かれたのではなく、局所 $U(1)$ ゲージ対称性の要請から定まっていることである。

\subsection{運動方程式}
このラグランジアンから、場の運動方程式が求まる。まず $\bar{\psi}$ で変分すると
\[
\left(i\gamma^\mu D_\mu - m\right)\psi = 0
\]
を得る。これは電磁場と相互作用するディラック方程式である。

次に $A_\mu$ で変分すると
\[
\partial_\nu F^{\nu\mu} = e\bar{\psi}\gamma^\mu\psi = e j^\mu
\]
を得る。これはマクスウェル方程式の場の理論版であり、右辺の電流 $j^\mu$ が電磁場の源になっている。すなわち、電子場が電磁場を作り、電磁場がまた電子場に作用する、という相互作用の相互性がこの二つの方程式に表れている。

\subsection{量子化すると何が相互作用するのか}
ここまでは古典場としての記述であった。これを量子化すると、ディラック場 $\psi$ の量子が電子・陽電子であり、電磁場 $A_\mu$ の量子が光子である。

すると相互作用項
\[
\mathcal{L}_{\mathrm{int}} = - e\bar{\psi}\gamma^\mu A_\mu \psi
\]
は、電子や陽電子が光子を放出したり吸収したりできることを表す。したがって、電子同士の電磁相互作用は、量子論的には光子のやり取りとして理解される。

この段階では、「異なる場の相互作用」とは何かがすでに明瞭である。すなわち、電子場だけ、あるいは電磁場だけを別々に考えるのではなく、ひとつのラグランジアンの中で両者を結び付け、その結び付きが量子化の後には粒子どうしの放出・吸収として現れるのである。

\subsection{QED の意義}
QED は、異なる場の相互作用を記述する最初の本格的な理論である。その本質は、電子場と電磁場を単に「結び付ける」のではなく、局所的な位相対称性が許す形で結び付けている点にある。

この理論は、電荷保存、電磁場の生成、電子・陽電子と光子の相互作用をひとつの枠組みで与える。さらに、「対称性が相互作用の形を決める」という考え方は、弱い相互作用や強い相互作用を記述する理論にもそのまま受け継がれている。

この意味で、場の量子論とは単に「相対論的な量子力学」ではなく、対称性に基づいて相互作用の形を決めるための統一的な言語なのである。
